% MT059 / explanation — canonical editable content.
从学习对象、边界形状、训练方法、扩展方式和参数数量五个方面比较两种模型。


\textbf{学习对象不同。} LR直接写出\(P(y=1\mid x)=\sigma(w^Tx+b)\)，从带标签数据学习\(w,b\)。Gaussian NB先学\(P(y=c)\)及\(p(x\mid c)\)，再用Bayes公式得到后验；因此前者判别式，后者生成式。


\textbf{共享方差时都可得到线性边界。} 看一维就能理解抵消过程。设两类共用方差\(v>0\)，均值为\(\mu_0,\mu_1\)，先验为\(\pi_0,\pi_1>0\)。取后验比是为了比较两类：两者相等时比为1、log为0，正好对应边界。对这个比取log，共同分母和高斯归一化系数消掉：
\[
\log\frac{P(1\mid x)}{P(0\mid x)}
=\log\frac{\pi_1}{\pi_0}-\frac{(x-\mu_1)^2-(x-\mu_0)^2}{2v}.
\]


把平方展开，分子里的\(x^2\)一正一负抵消，得到
\[
\log\frac{P(1\mid x)}{P(0\mid x)}
=\frac{\mu_1-\mu_0}{v}x+\frac{\mu_0^2-\mu_1^2}{2v}+\log\frac{\pi_1}{\pi_0}.
\]


这正是\(wx+b\)的形式。多维NB把各维对应项相加，得到\(w^Tx+b\)；若每类方差不同，平方项通常不会抵消。记\(p=P(1\mid x)\)，则另一类的概率为\(1-p\)。上式说明\(\log[p/(1-p)]=f\)，即\(p/(1-p)=e^f\)。两边乘以\(1-p\)并把含\(p\)的项移到一起，得\(p(1+e^f)=e^f\)，所以\(p=1/(1+e^{-f})\)。这就是sigmoid。


\textbf{训练原则可都叫MLE，但最大化的量不同。} LR让给定输入后的真实标签更可信；NB让观察到的输入与类别在联合模型下更可信。NB的高斯参数可用分类后的均值、方差等估计；普通LR通常需要迭代数值优化。二者都可扩展多分类，也都能加入正则或参数先验。


\textbf{数参数前说明“共享”的含义。} 二分类\(d\)维LR有\(d\)个权重和1个偏置，共\(d+1\)。配套答案采用\(\Sigma_0=\Sigma_1=\sigma^2I\)：两类均值共\(2d\)个数，所有维度共用1个方差，类条件部分共\(2d+1\)。若还学习先验\(\pi\)，另一个先验由\(1-\pi\)决定，只再加1。


补充\(d=2\)的清单：\((\mu_{01},\mu_{02},\mu_{11},\mu_{12},\sigma^2)\)是5个数，加\(\pi\)为6个。若仅要求同一维在类间共享、不同维各有方差\(v_1,v_2\)，则是4个均值＋2个方差＋1个先验，共7个。英文作答采用原解的共享标量约定，并注明是否计先验，不能混用两种计数。


\textbf{出处：} 2023B期中 Q8，原卷第5页（10分）


\textbf{答案依据：} 2023B随卷解答 Q8随卷答案；补清2d+1的模型约定和先验计数。
